% !TeX root = ../navier-stokes-manuscript.tex
% ============================================================
% 2.1 Vortex Stretching
% ============================================================

The momentum equation has four pieces:
\[
  \underbrace{\partial_t u}_{\text{local acceleration}}
  +
  \underbrace{(u\cdot\nabla)u}_{\text{nonlinear transport}}
  =
  \underbrace{-\nabla p}_{\text{pressure}}
  +
  \underbrace{\nu\Delta u}_{\text{viscous diffusion}}.
\]
The central regularity problem lies in the competition between nonlinear transport and viscosity.
The hope is simple: viscosity smooths the fluid faster than nonlinear transport can sharpen it.
If that remains true at every scale, the velocity stays smooth.

\subsection{Vortex Stretching}\label{sub:ThePerfectlyStretchingVortex}

In a narrow material vortex tube, the surrounding flow can pull a material segment of the axis apart while the same rotating fluid keeps moving through time.
The part of the velocity gradient that records this deformation is the symmetric strain
\[
  S:=\frac12\bigl(\nabla u+(\nabla u)^{\mathsf T}\bigr),
\]
and if \(e\) points along the segment while \(L(t)\) denotes its length, extension along \(e\) at rate \(\sigma(t)>0\) means
\[
  Se=\sigma(t)e,
  \qquad
  \sigma(t)>0,
  \qquad
  \frac{\dot L(t)}{L(t)}
  =
  e\cdot Se
  =
  \sigma(t).
\]
The same strain that separates the ends is lengthening that one material piece along its own axis.

Because the fluid is incompressible, the same motion squeezes the cross-section.
Let \(\mathcal V\) be the fixed material volume, set \(A(t):=\mathcal V/L(t)\), and write \(r_{\mathrm{eff}}(t):=\sqrt{A(t)/\pi}\).
Then
\[
  \nabla\cdot u=0,
  \qquad
  \frac{d}{dt}(A L)=0,
  \qquad
  \frac{\dot A}{A}=-\sigma(t),
  \qquad
  \frac{\dot r_{\mathrm{eff}}}{r_{\mathrm{eff}}}
  =
  -\frac{\sigma(t)}{2}.
\]
Lengthening and narrowing are one coupled event: as the axis is drawn out, the core radius falls.

That same deformation acts on the tube's rotation.
With \(\omega:=\nabla\times u\), the vorticity satisfies along the moving fluid
\[
  \frac{D\omega}{Dt}
  =
  S\omega+\nu\Delta\omega,
  \qquad
  \frac{D}{Dt}
  :=
  \partial_t+u\cdot\nabla,
\]
and when the vorticity aligns with the stretching direction, \(\omega=|\omega|e\),
\[
  S\omega
  =
  \sigma(t)\omega,
  \qquad
  \frac12\frac{D|\omega|^2}{Dt}
  =
  \sigma(t)|\omega|^2
  +
  \nu\,\omega\cdot\Delta\omega.
\]
Axial stretching now appears twice inside the same material segment: it lengthens the tube and amplifies aligned vorticity.
The viscous term remains in that balance, with no pointwise sign forced by the identity itself, so growth occurs only on intervals where the full right-hand side stays positive.

What is growing here is the rotating part of the velocity gradient, not the kinetic energy of the segment as a whole.
If the speed stays bounded by \(U_0\) on the fixed-volume material segment \(\Omega_{\mathrm{seg}}(t)\), then
\[
  E_{\mathrm{seg}}(t)
  :=
  \frac12\int_{\Omega_{\mathrm{seg}}(t)}|u(x,t)|^2\,dx
  \leq
  \frac12U_0^2\mathcal V
  <
  \infty,
  \qquad
  \left|\frac12\bigl(\nabla u-(\nabla u)^{\mathsf T}\bigr)\right|_{\mathrm F}
  =
  \frac{|\omega|}{\sqrt2}
  \leq
  |\nabla u|_{\mathrm F}.
\]
A bounded-speed segment of fixed material volume can keep finite kinetic energy while its rotating gradient grows and its core narrows.
The gradients rise.

\divider{\NavierStokes{} Scaling}

Once the core has narrowed to width \(\ell\), the equation has to be read at that width.
The comparison comes from the \NavierStokes{} rescaling, not from pretending that a later state of this same tube has already been reached.
Fix \(t_0\) and set \(\ell:=r_{\mathrm{eff}}(t_0)\).
For \(\lambda>1\), define
\[
  u_\lambda(x,t)
  :=
  \lambda u(\lambda x,\lambda^2t),
  \qquad
  p_\lambda(x,t)
  :=
  \lambda^2p(\lambda x,\lambda^2t).
\]
At time \(\lambda^{-2}t_0\), the corresponding vortex in \(u_\lambda\) has width \(\lambda^{-1}\ell\).
Differentiating the rescaled fields gives
\[
\begin{aligned}
  \partial_tu_\lambda(x,t)
  &=
  \lambda^3(\partial_tu)(\lambda x,\lambda^2t),
  \\
  (u_\lambda\cdot\nabla)u_\lambda(x,t)
  &=
  \lambda^3\bigl((u\cdot\nabla)u\bigr)(\lambda x,\lambda^2t),
  \\
  \nabla p_\lambda(x,t)
  &=
  \lambda^3(\nabla p)(\lambda x,\lambda^2t),
  \\
  \nu\Delta u_\lambda(x,t)
  &=
  \lambda^3\nu(\Delta u)(\lambda x,\lambda^2t).
\end{aligned}
\]
The narrowed width changes none of the equation's internal proportions.
Local acceleration, nonlinear transport, pressure, and viscosity all acquire the same \(\lambda^3\) factor.
The sharpening term has not lost ground to viscosity by moving to the finer scale.

\divider{Critical Height}

So the balance survives the change of width.
What can still change is the size used to record that balance, because a norm may grow or shrink under rescaling even when the equation itself stays level.
Let \(\Lambda:=(-\Delta)^{1/2}\).
At Sobolev height \(s\), a change of variables gives
\[
  \norm{u_\lambda(t)}_{\dot H^s}^2
  =
  \lambda^{2s-1}
  \norm{u(\lambda^2t)}_{\dot H^s}^2.
\]
The factor \(\lambda^{2s-1}\) is contributed by scaling alone.
It disappears only at \(s=\tfrac12\).
Set
\[
  H(t)
  :=
  \frac12\norm{u(t)}_{\dot H^{1/2}}^2
  =
  \frac12\norm{\Lambda^{1/2}u(t)}_2^2.
\]
At this height the scale change itself is silent, while the levels below and above it still move:
\[
  \norm{u_\lambda(t)}_2^2
  =
  \lambda^{-1}\norm{u(\lambda^2t)}_2^2,
  \qquad
  H_\lambda(t)=H(\lambda^2t),
  \qquad
  \norm{\nabla u_\lambda(t)}_2^2
  =
  \lambda\norm{\nabla u(\lambda^2t)}_2^2.
\]
Kinetic energy falls under concentration, first-derivative energy rises, and the critical height stays fixed.
Finite energy still leaves the concentration problem open.
With the ruler now fixed, any rise or fall in \(H\) has to come from the equation itself.

\divider{Nonlinear Production versus Viscous Dissipation}

Apply \(\Lambda^{1/2}\) to the momentum equation and pair with \(\Lambda^{1/2}u\).
Write the signed contribution of nonlinear transport as
\[
  \mathcal P_{1/2}(t)
  :=
  -\operatorname{Re}
  \pair
    {\Lambda^{1/2}u(t)}
    {\Lambda^{1/2}\bigl((u(t)\cdot\nabla)u(t)\bigr)}_{L^2}.
\]
The pressure pairing vanishes because \(\nabla\cdot u=0\), while the Laplacian contributes \(-\nu\norm{\Lambda^{3/2}u}_2^2\).
The two effects meet in the exact evolution law
\[
  H'(t)
  +
  \nu\norm{\Lambda^{3/2}u(t)}_2^2
  =
  \mathcal P_{1/2}(t).
\]
The critical height rises only when nonlinear production exceeds viscous dissipation.
Under the same rescaling,
\[
  \mathcal P_{1/2}[u_\lambda](t)
  =
  \lambda^2\mathcal P_{1/2}[u](\lambda^2t),
  \qquad
  \nu\norm{\Lambda^{3/2}u_\lambda(t)}_2^2
  =
  \lambda^2\nu\norm{\Lambda^{3/2}u(\lambda^2t)}_2^2.
\]
Across the comparison family, neither contribution gains any relative strength.
Both carry the same \(\lambda^2\) factor, and that common factor comes with the contracted time coordinate \(t\mapsto\lambda^{-2}t\).
The spatial balance is still unresolved, but it is now being read on a shorter clock.

\divider{The Heat Clock}

Viscosity already carries that clock inside its linear action.
For the heat equation \(v_t=\nu\Delta v\),
\[
  \widehat v(\xi,t+\delta t)
  =
  e^{-\nu|\xi|^2\delta t}\widehat v(\xi,t).
\]
A structure localized to width \(r\) lives at frequencies \(|\xi|\simeq r^{-1}\).
An order-one viscous change appears when \(\nu|\xi|^2\delta t\simeq1\), so viscosity associates that width with the time
\[
  \tau_\nu(r)
  :=
  \frac{r^2}{\nu}.
\]
This is not the declared lifetime of the material vortex.
It is the time viscosity alone needs to alter a structure of width \(r\) by order one.
At the finer comparison width,
\[
  \tau_\nu(\lambda^{-1}r)
  =
  \lambda^{-2}\tau_\nu(r),
\]
so the viscous comparison time contracts by the same factor as the full \NavierStokes{} time rescaling.
Every finer spatial test is paired with a quadratically shorter viscous clock, while the nonlinear and viscous contributions remain evenly matched.

\divider{The Zeno Cascade}

Repeated halving makes that pairing explicit.
For
\[
  r_n:=2^{-n}r_0,
  \qquad
  \tau_n:=\frac{r_n^2}{\nu},
\]
the successive clocks satisfy
\[
  \tau_n
  =
  \frac{r_0^2}{\nu}4^{-n},
  \qquad
  \sum_{n=0}^{\infty}\tau_n
  =
  \frac{4r_0^2}{3\nu}
  <
  \infty.
\]
Infinitely many finer comparison scales can fit inside finite total comparison time.
But those scales are only a schedule until one actual solution realizes them.
A candidate singular history would have to pass through widths
\[
  r_0>r_1>r_2>\cdots\longrightarrow0
\]
at times
\[
  t_0<t_1<t_2<\cdots\uparrow T_*<\infty,
  \qquad
  t_{n+1}-t_n\asymp\tau_n,
\]
with comparison constants independent of \(n\).
The remaining time would then satisfy
\[
  T_*-t_n
  \asymp
  \sum_{k=n}^{\infty}\tau_k
  =
  \frac{4r_n^2}{3\nu}.
\]
By the stage where the width has reached \(r_n\), only order \(r_n^2/\nu\) of comparison time remains.
So the unresolved competition does not disappear at finer scale.
It returns as a same-history demand: one smooth velocity field would have to keep producing narrower concentration through infinitely many scales while the available viscous clock collapses quadratically toward zero.
